% Section "Multireference geometries: attempted, out of scope"
% (sec:mr-whynotmr). Gehoert in das Outlook-/Future-work-Kapitel der
% Thesis, nicht in das OMol25-Kapitel; dort steht nur ein Verweis
% (Methods, DFT Reevaluation). Lange Vorfassung: section_mr_feasibility.tex.

\section{Multireference geometries: attempted, out of scope}
\label{sec:mr-whynotmr}

Multireference geometries and energies were attempted for these reactions
because they would be the gold standard for reactions with high
multireference character. They were placed out of scope. Problems arose,
and they are not needed for the question of this thesis: the models were
trained on \mbox{$\omega$B97M-V/def2-TZVPD} and are trying to find the
transition states of that surface, so the reference stays at that level.
A reference closer to the exact surface would answer a different
question, and a model agreeing with it would be right for the wrong
reason. No multireference barrier enters this work.

\paragraph{The overall problem: active space selection.} Every
multireference method needs an active space, and the result depends on
it. autoCAS~\cite{stein2019autocas} selects the space from a DMRG
calculation over all valence orbitals, which was not practical for
molecules of this size. AVAS~\cite{sayfutyarova2017} and selection by
hand require experience that could not be acquired within the time of
this work. A smaller choice of molecules for this part of the test set
would have avoided this. This is also why the standard protocol, NEVPT2
energies on UKS geometries, was not carried out.

\paragraph{The additional problem for geometries: active space
consistency.} Starting from the DFT transition state, AVAS selected a
space, CASSCF optimised a transition state with it, and AVAS at the
optimised geometry returned a different space. The geometry had been
optimised on a surface that does not belong to it. A possible solution is
an iteration:

\begin{enumerate}
  \item select the active space, with autoCAS where affordable;
  \item optimise the transition state at the CASPT2 or NEVPT2 level;
  \item re-select the active space at the optimised geometry;
  \item repeat until the selected space is the one the geometry was
    optimised with;
  \item add the next orbitals to the converged space and check that the
    result does not change.
\end{enumerate}

The re-selection in the optimised orbital basis is what the autoCAS
authors recommend for a single structure~\cite{stein2019autocas}; the
iteration above applies it to a moving geometry.

Significant time of this work went into this attempt. To the author's
knowledge, transition-state geometries optimised at a multireference
level do not exist as a benchmark set, which is why it was worth trying
and is documented here.

